\documentclass[12pt,a4paper]{article}
\usepackage[margin=2cm]{geometry}
\usepackage{graphicx}
\usepackage{enumitem}
\usepackage{hyperref}
\usepackage{fancyhdr}
\usepackage{titlesec}
\usepackage{xcolor}
\usepackage{float}

\pagestyle{fancy}
\fancyhf{}
\rhead{Drone Mapper -- High Level Design}
\lhead{Advanced Topics in Programming 2026B}
\cfoot{\thepage}

\titleformat{\section}{\Large\bfseries}{}{0em}{}
\titleformat{\subsection}{\large\bfseries}{}{0em}{}
\titlespacing*{\section}{0pt}{8pt}{4pt}
\titlespacing*{\subsection}{0pt}{6pt}{3pt}

\begin{document}

\begin{center}
{\LARGE\bfseries High Level Design Document}\\[4pt]
{\Large Drone Mapper -- Exercise 1}\\[8pt]
{\large Advanced Topics in Programming, Semester B 2026}\\[12pt]
\begin{tabular}{ll}
\textbf{Name} & \textbf{ID} \\
\hline
Almog Tavor & (ID: 323084962) \\
Yonatan Kahana & (ID: 212223036) \\
\end{tabular}
\end{center}

\vspace{6pt}

% ============================================================
\section{1. UML Class Diagram}
% ============================================================

The diagram below shows the main classes, their relationships, and key
methods. Interfaces are blue, core classes yellow, mock implementations
green, data types pink, and utilities purple.

\vspace{2pt}
\begin{center}
\includegraphics[width=\textwidth,height=0.72\textheight,keepaspectratio]{class_diagram.pdf}
\end{center}

\vspace{-4pt}
\noindent Key relationships:
\begin{itemize}[nosep,topsep=2pt]
\item \textbf{Simulator} owns all mocks, the Drone, both maps, and orchestrates the loop.
\item \textbf{Drone} sees only abstract interfaces (ILidarSensor, IPositionSensor,
      IMovementDriver) through the command/response pattern -- it never touches the
      mocks or truth directly.
\item \textbf{MockWorld} is shared state between the three mocks so that movement
      updates are visible to the position and lidar sensors.
\item \textbf{VoxelGrid} is the shared spatial data structure, composed into both
      BuildingTruth and BuildingMap.
\end{itemize}

\newpage
% ============================================================
\section{2. UML Sequence Diagram -- Main Flow}
% ============================================================

The diagram shows the four phases of the main simulation loop: location
query, scanning, movement, and scoring.

\vspace{2pt}
\begin{center}
\includegraphics[width=\textwidth,height=0.85\textheight,keepaspectratio]{sequence_diagram.pdf}
\end{center}

\newpage
% ============================================================
\section{3. Design Considerations and Alternatives}
% ============================================================

\subsection{3.1 Separation of Drone Logic from Simulation}

The drone never directly accesses the building truth or the mock implementations.
Instead, the \textbf{Simulator} acts as a mediator: it asks the Drone for a
command, dispatches it to the appropriate mock, and feeds the result back. This
mirrors a real embedded system where the flight controller issues commands and
receives sensor data through hardware abstraction layers.

\textbf{Alternative considered:} Letting the drone call sensor/driver interfaces
directly (dependency injection). We rejected this because the assignment's ``main
loop'' specification explicitly requires the simulator to ask the drone for
commands, making the command-response pattern the natural fit.

\subsection{3.2 Voxel Grid as Shared Spatial Representation}

Both \texttt{BuildingTruth} and \texttt{BuildingMap} wrap the same
\texttt{VoxelGrid} class. This ensures identical coordinate-to-cell mapping
and makes the scoring comparison straightforward (iterate both grids in lock-step).

\textbf{Alternative considered:} Using an octree or sparse hash map for the
drone's known map (since most cells start as ``unmapped''). We chose a dense
grid for simplicity and because the building sizes in exercise~1 are small
enough that memory is not a concern.

\subsection{3.3 Full Lidar Frame Back-Projection}

Initially we used only the central ray of each cardinal scan to update the map.
This left corner cells unmapped and produced scores around 96\%. We improved
to back-projecting \emph{every} ray in the lidar frame to world coordinates,
walking each ray through the voxel grid and marking cells accordingly. This
raised scores to 99\%+ and reduced command counts by 60\%.

\textbf{Trade-off:} The full frame approach is $O(n^2)$ per scan (where $n$
is the lidar grid side), but with our capped grid size ($\le 51 \times 51$)
and only 6 scans per waypoint, the overhead is negligible.

\subsection{3.4 Frontier-Based BFS Exploration}

The drone maintains a ``frontier'' -- known-empty cells adjacent to unmapped
cells. At each waypoint it runs BFS through known-empty cells to find the
nearest frontier, plans the shortest path, and follows it. This is complete
(it will visit every reachable cell) and deterministic (BFS neighbor order
is fixed).

\textbf{Alternative considered:} DFS / wall-following. DFS is simpler but can
lead to long back-tracking paths. A\textsuperscript{*} with distance heuristic
could reduce path lengths but adds complexity without material benefit at
exercise~1 grid sizes.

\subsection{3.5 Strong Types Without mp-units}

The assignment recommends mp-units, but it is a heavy external dependency.
We provide a lightweight \texttt{Units.h} header with \texttt{Length} and
\texttt{Angle} classes that support the same multiplication syntax
(\texttt{5 * cm}, \texttt{90 * deg}). This keeps the build self-contained
while enforcing unit safety at the type level.

\subsection{3.6 Polygon Boundary Support}

The mission boundary is defined as an arbitrary polygon (supporting concave
shapes like L-buildings). We use the standard ray-casting point-in-polygon
algorithm to classify each voxel during \texttt{BuildingMap} initialization.

\textbf{Alternative considered:} Axis-aligned bounding box (simpler, but cannot
represent non-rectangular buildings). We chose polygon support since the
assignment explicitly mentions ``polygon of \{X, Y\}''.

\subsection{3.7 Floating-Point Position Snapping}

We discovered that cumulative floating-point drift from \texttt{cos($\pi$/2)}
not being exactly zero caused the drone's actual position to cross cell
boundaries, leading to incorrect lidar readings and infinite loops. The fix:
snap world position to the nearest 0.01\,cm after every movement in the mock.

\subsection{3.8 Error Recovery Strategy}

All input file errors are recoverable where possible: missing config keys use
defaults, malformed map rows are padded, and unknown keys are logged. Only truly
fatal errors (file not found, no grid data) cause the program to print an error
and exit via \texttt{main} return -- never via \texttt{exit()} or exceptions
escaping \texttt{main}.

% ============================================================
\section{4. Testing Approach}
% ============================================================

\begin{itemize}[nosep]
\item \textbf{Bottom-up:} Each module is tested in isolation before integration.
      Mocks are tested against known truth grids with predictable geometry.
\item \textbf{Integration:} Full truth$\to$simulator$\to$drone pipeline tests
      verify end-to-end behaviour (score threshold, no collisions).
\item \textbf{Error paths:} Missing files, unknown keys, and malformed values
      are verified to be handled gracefully with recoverable defaults.
\item \textbf{Determinism:} All tests are deterministic. Same input always
      produces the same output.
\end{itemize}

\subsection{4.1 Sample Input Sets}

Three input sets exercise different scenarios:

\begin{enumerate}[nosep]
\item \textbf{set1} -- Simple 10$\times$10 empty room (baseline correctness).
\item \textbf{set2} -- 15$\times$12 building with a sealed 3$\times$3 pocket
      (tests that inaccessible areas remain unmapped without crashing).
\item \textbf{set3} -- L-shaped building with non-convex polygon boundary
      (tests polygon clipping and multi-arm exploration).
\end{enumerate}

\end{document}
